% =====================================================================
% TRACK B -- DETAILED METHODS & RESULTS  (for the Overleaf appendix)
%
% To append to your Overleaf project:
%   - copy from the intro paragraph (right after \maketitle) through the
%     last \section, i.e. all the \section blocks + table floats. Do NOT
%     copy \documentclass / \title / \maketitle. Drop it under \appendix.
%   - your preamble must load: amsmath, amssymb, booktabs, array, xcolor.
%   The body uses NO custom macros and \label/\ref, so it renumbers
%   automatically inside your paper -- copy-paste safe.
% This file also compiles standalone (pdflatex x2 for refs) for preview.
% =====================================================================
\documentclass[11pt,a4paper]{article}
\usepackage[margin=2.3cm]{geometry}
\usepackage[T1]{fontenc}
\usepackage{lmodern}
\usepackage{microtype}
\usepackage{amsmath,amssymb}
\usepackage{booktabs}
\usepackage{array}
\usepackage{xcolor}
\usepackage[colorlinks=true,linkcolor=blue!55!black,citecolor=blue!55!black]{hyperref}
\renewcommand{\arraystretch}{1.15}
\title{\vspace{-1.4em}\textbf{Track B: Detailed Methods and Results}}
\author{}
\date{2026-07-30}

\begin{document}
\maketitle
\vspace{-2em}

\noindent This document expands the compact Track~B handoff into a full methods-and-results reference suitable
for a paper appendix: it defines every method and metric, states the exact computations, and reports all
result tables gathered so far. Track~B is the set of eight follow-up experiments (B1--B8) requested in the TMLR
revision to strengthen the TraLO paper. \emph{TraLO} trains a network to satisfy a \textbf{transductive
prediction-count constraint} (an upper bound on how many test points may be predicted as a designated
``constrained'' class) by adding a soft, bounded penalty to the training loss, and is compared against
constrained-optimization duals, post-hoc clipping, and imbalanced-learning baselines.

\medskip
\noindent\textbf{Status.} B1--B8 are complete except B1's third backbone (ViT-B/16), which is still computing;
its rows are marked \emph{pending} in Table~\ref{tab:b1} and do not affect any stated conclusion.

% =====================================================================
\section{Problem setup and notation}
% =====================================================================
We are given a trained-on set and an unlabeled \emph{test} (transductive) set. One class $c^\star$ is
designated \emph{constrained}. A deployment rule caps the number of test points predicted as $c^\star$:
\begin{equation}
\underbrace{\textstyle\sum_{i\in\text{test}} \mathbf{1}[\hat y_i = c^\star]}_{\text{hard count}} \;\le\; K,
\qquad K = (\text{cap fraction})\times(\text{true prevalence of }c^\star).
\end{equation}
The \textbf{soft count} $S_{c^\star}=\sum_i p_i(c^\star)$ (sum of softmax probabilities) is the differentiable
surrogate used inside the loss; the \textbf{hard count} (using $\arg\max$) is used only to \emph{verify}
satisfaction. The \emph{excess} is $E=\max(0,\,S_{c^\star}-K)$. Caps come in two flavours: a \emph{global} cap
$G$ on the total count and \emph{local} caps $L$ on per-group counts; a cell tagged \texttt{L30\_G30} sets both
to $30\%$ of the true prevalence (symmetric). Smaller percentages are \emph{tighter} caps and bind harder.

% =====================================================================
\section{Methods compared}
% =====================================================================
All methods produce a classifier and then must meet the cap. They differ in \emph{how}: TraLO and the duals
\emph{train} to satisfy it; the imbalanced baselines train an imbalance-aware classifier and then clip;
post-hoc clipping trains normally and only edits predictions at inference.

\subsection{TraLO (proposed)}
Two phases. \textbf{(i) Warm-up:} standard cross-entropy (CE) for $W$ epochs. \textbf{(ii) Constraint phase}
(up to $300$ epochs): minimise
\begin{equation}
\mathcal{L} \;=\; \mathcal{L}_{\text{CE}} \;+\; \lambda_g\,\mathcal{L}_{\text{global}} \;+\;
\lambda_\ell\,\mathcal{L}_{\text{local}},\qquad
\mathcal{L}_{\text{pen}} = \underbrace{\frac{E}{E+K}}_{\text{rational saturation}} \;+\;
\underbrace{\rho\,\frac{(E/K)^2}{1+(E/K)^2}}_{\text{bounded quadratic}} \;\in\; [0,\,1+\rho).
\label{eq:pen}
\end{equation}
The penalty is deliberately \textbf{bounded} (it saturates as the violation grows), unlike the duals' linear
penalty. Training uses several controls, each individually ablated in the paper: a \emph{$\lambda$-ratchet}
(increment $\lambda$ each epoch until the cap is met, then freeze); a \emph{linear $\rho$ schedule} (raise
$\rho$ until first satisfaction, then freeze); a \emph{hinge / $\lambda$-toggle} (zero the multipliers when
satisfied, restore when violated, with oscillation detection); an \emph{optimizer reset at first satisfaction}
(reset the Adam moment estimates to kill post-satisfaction momentum); a \emph{CE-saturation skip} (end phase~ii
when train accuracy $\ge 0.995$); and \emph{convergence early-stopping} (5 consecutive satisfied epochs).

\subsection{Constrained-optimization duals}
These solve the constrained problem with a Lagrange multiplier $\lambda$ updated by dual ascent.
\begin{itemize}
\item \textbf{Fioretto-LDF} (linear dual): penalty $\lambda(S_{c^\star}-K)$ with
$\lambda \leftarrow [\lambda + \eta\,(S_{c^\star}-K)]_+$. The linear penalty can \emph{wind up} ($\lambda$
grows without bound if the cap is never met).
\item \textbf{ALM} (augmented Lagrangian): $\lambda \leftarrow \max(0,\lambda+\eta(S_{c^\star}-K)) +
\mu\,(S_{c^\star}-K)^+$ with $\mu$ growing linearly. This is the standard literature fix for linear-penalty
windup, and empirically a \emph{stronger} baseline than plain Fioretto.
\item \textbf{Hounie-RCL} (resilient constrained learning): allows the constraint level to relax adaptively at
a cost. In the tight-cap regime it underperforms sharply (see B8).
\end{itemize}

\subsection{Post-hoc clipping (Shifman-LP)}
Train the network with plain CE (no constraint), then at inference solve a small linear program that flips the
minimal set of borderline predictions so the hard count meets $K$. It is \textbf{post-hoc}: it never trains to
satisfy the cap and therefore \emph{cannot re-rank} the decision boundary; it only relabels the least-confident
predictions. It is a different category of method from the trained ones, and is reported separately.

\subsection{Imbalanced-learning baselines (B1)}
These target the rare constrained class directly. Each trains the backbone with an imbalance-aware loss during
the warm-up phase and then applies Shifman-LP clipping; the warm-up is therefore their \emph{entire} training
budget. Let $p_y$ be the softmax probability of the true class $y$, $n_c$ the number of training samples in
class $c$, and $\pi_c$ the class prior (frequency).
\begin{itemize}
\item \textbf{Focal loss} (Lin et al., 2017): $\ \ell_{\text{focal}} = -\,\alpha\,(1-p_y)^{\gamma}\,\log p_y$,
with $\alpha=0.25,\ \gamma=2.0$. The modulating factor $(1-p_y)^\gamma$ down-weights easy (high-confidence)
examples, concentrating gradient on hard / rare cases and boosting minority-class recall.
\item \textbf{Class-balanced weighting} (Cui et al., 2019): weight class $c$'s loss by the inverse
\emph{effective number} of samples, $\ w_c = \dfrac{1-\beta}{1-\beta^{\,n_c}}$, with $\beta=0.9999$. The
effective number $(1-\beta^{n_c})/(1-\beta)$ discounts information overlap among many similar samples, so rare
classes receive proportionally larger weight.
\item \textbf{Logit adjustment} (Menon et al., 2020): add a prior-based offset to the logits,
$\ z_c \leftarrow z_c + \tau\,\log \pi_c$, with $\tau=1.0$. This shifts the decision threshold to compensate
for label imbalance and is Bayes-consistent for the balanced error rate.
\end{itemize}

% =====================================================================
\section{Metrics, columns, and parameters}
% =====================================================================
\begin{itemize}
\item \textbf{cc-F1 (constrained-class F1)} -- the F1 score of the constrained class $c^\star$ alone; the
primary deployment metric. Because a cap at fraction $c$ of the true count limits recall to $\le c$, the
attainable cc-F1 is \emph{ceiling-bounded} by $2c/(1+c)$ (e.g.\ $c=0.3 \Rightarrow 0.46$). Absolute cc-F1 gaps
are therefore small by construction (``ceiling crush'').
\item \textbf{macro-F1} -- the unweighted mean of per-class F1; overall classification quality across all
classes. It is the metric the imbalanced baselines (B1) target.
\item \textbf{Native satisfaction} -- the fraction of runs whose hard (argmax) count already meets the cap at
inference \emph{without} any post-hoc flip. Methods that train to satisfy (TraLO, duals) reach $1.0$; methods
that clip afterwards (post-hoc clip, imbalanced baselines) are $0.0$ and rely on the LP.
\item \textbf{Cap tag} \texttt{L$X$\_G$X$} -- local and global caps each set to $X\%$ of the true prevalence of
$c^\star$; smaller $X$ = tighter.
\item \textbf{Warm-up epochs $W$} -- CE-only pre-training before the constraint phase, and the ``headroom''
lever: short $W$ leaves CE unsaturated (train-acc $<1$) so the penalty can re-rank the boundary; long $W$
saturates CE (train-acc $\to 1$), leaving only a uniform logit shift.
\item \textbf{Gap ``$A-B$''} -- the matched-seed mean of (metric$_A$ $-$ metric$_B$) within a cell; positive
means $A$ is better.
\item \textbf{Cell} -- the atomic unit of analysis: one (dataset, backbone, cap, method) averaged over the 4
seeds. Backbones, datasets, and caps are \textbf{never pooled}.
\item \textbf{W/T/L at $\tau$} -- across cells, the counts with gap $\ge +\tau$ (win), $|\text{gap}|<\tau$
(tie), and $\le -\tau$ (loss).
\end{itemize}

% =====================================================================
\section{Statistical methodology (the computations)}
% =====================================================================
\begin{itemize}
\item \textbf{Matched-seed (paired) comparison.} A seed fixes the initialization, data order, and warm-up, so
two methods at the same seed start from the same network; we compare their per-seed difference. This pairing
removes between-seed variance, so we \emph{never} use pooled-standard-deviation (unpaired) tests.
\item \textbf{Bootstrap confidence intervals.} For a comparison at a given cap, we resample the constituent
cells (backbones) with replacement $20{,}000$ times and take the $2.5$/$97.5$ percentiles of the resampled mean
gap. A CI that excludes $0$ is significant.
\item \textbf{Benjamini--Hochberg FDR.} Across the family of per-cell paired tests we control the false
discovery rate at $q=0.05$ and report how many cells survive.
\item \textbf{Win-bar rule.} A regime counts as a ``win'' if the mean paired cc-F1 gap $\ge \tau$ \emph{and} at
least half the cells individually clear $\tau$; we sweep $\tau\in\{0.003,0.005,0.010\}$ to show the
classification is not threshold-sensitive.
\end{itemize}
Every experiment ran $4$ seeds ($1$--$4$); the campaign totals $568$ runs (plus $60$ ViT-B/16 runs for B1 in
progress), with $0$ failures.

% =====================================================================
\section{Experiments and results}
% =====================================================================

\subsection{B1 -- Imbalanced-learning baselines}
\textbf{Question.} Does an imbalance-aware clipper (focal / class-balanced / logit-adjust $+$ clipping) close
the macro-F1 gap the paper attributes to constraint-time training? \textbf{What ran.} $252$ runs: the $5$
methods $\times$ \{Oct, Derm, Tissue\}MNIST $\times$ \{MobileNetV3, RegNetY-400MF\} $\times$ warm-up
$\{1,5,50\}$ $\times$ cap L30 $\times$ 4 seeds. Table~\ref{tab:b1} reports the macro-F1 gap
(TraLO $-$ baseline) at warm-up 50. \textbf{Found.} Only focal is competitive; it \emph{beats} TraLO on
Derm/MNV3 and \emph{loses} on Oct/MNV3, i.e.\ the outcome is backbone- and dataset-dependent. cc-F1 is tied
(ceiling). The durable, universal separation is native satisfaction ($1.0$ vs $0.0$). \emph{Verdict: an honest
tie on macro-F1, plus native satisfaction.}

\begin{table}[h]\centering\small
\begin{tabular}{llrrrr}
\toprule
dataset & backbone & vs focal & vs class-bal & vs logit-adj & vs clip \\
\midrule
Derm   & MobileNetV3 & $-0.012$ & $-0.016$ & $+0.021$ & $+0.010$ \\
Derm   & RegNetY400MF & $+0.006$ & $+0.022$ & $+0.042$ & $+0.013$ \\
Oct    & MobileNetV3 & $+0.025$ & $+0.023$ & $+0.021$ & $+0.023$ \\
Oct    & RegNetY400MF & $-0.005$ & $+0.022$ & $-0.000$ & $+0.022$ \\
Tissue & MobileNetV3 & $+0.020$ & $+0.014$ & $+0.026$ & $+0.029$ \\
Tissue & RegNetY400MF & $-0.002$ & $+0.003$ & $-0.001$ & $+0.001$ \\
\midrule
\multicolumn{6}{l}{\itshape ViT-B/16 (third backbone): runs in progress; rows to be appended.} \\
\bottomrule
\end{tabular}
\caption{B1: macro-F1 gap (TraLO $-$ baseline) at warm-up 50, cap L30, mean over 4 seeds. Positive $=$ TraLO
better. Focal is the only competitive baseline and is backbone-dependent.}
\label{tab:b1}
\end{table}

\subsection{B2 -- Native-resolution replication}
\textbf{Question.} Does the tight-cap behaviour reproduce at native ($224$px) resolution? \textbf{What ran.}
$192$ runs on five native MedMNIST datasets (Derm$=$HAM10000, Retina, Blood, OrganA, TissueNative) $\times$
warm-up $\{1,5\}$ $\times$ cap L30 ($+$L20) $\times$ 4 seeds. We compare against the best \emph{trained}
baseline (Fioretto-LDF); post-hoc clip is excluded here because at warm-up $1/5$ it is undertrained, so a gap
over it would measure ``clip did not train,'' not a real advantage; Hounie was not run natively.
\textbf{Found} (Table~\ref{tab:b2}): TraLO \emph{ties} the best trained baseline on all five datasets and both
warm-ups ($|\Delta|\le 0.011$) -- an honest negative replication, consistent with the mechanism (both trained
methods re-rank equally when CE has headroom). Native satisfaction remains $1.0$.

\begin{table}[h]\centering\small
\begin{tabular}{lccc}
\toprule
dataset & TraLO (w1 / w5) & best baseline, Fioretto (w1 / w5) & $\Delta =$ TraLO $-$ best (w1 / w5) \\
\midrule
Derm      & 0.331 / 0.412 & 0.325 / 0.412 & $+0.006$ / $-0.000$ \\
Retina    & 0.181 / 0.208 & 0.181 / 0.200 & $+0.000$ / $+0.008$ \\
TissueNat & 0.212 / 0.265 & 0.211 / 0.254 & $+0.001$ / $+0.011$ \\
OrganA    & 0.395 / 0.463 & 0.392 / 0.463 & $+0.003$ / $+0.000$ \\
Blood     & 0.460 / 0.457 & 0.460 / 0.457 & $+0.000$ / $+0.000$ \\
\bottomrule
\end{tabular}
\caption{B2: cc-F1 at native $224$px, TraLO vs the best trained baseline (Fioretto-LDF), each cell warm-up
$1$/$5$, mean over 4 seeds. A tie across all datasets and warm-ups.}
\label{tab:b2}
\end{table}

\subsection{B3 -- ALM (augmented-Lagrangian) baseline}
\textbf{Question.} Does ALM, the standard fix for linear-penalty windup, close the gap? \textbf{What ran.} $24$
runs (Fioretto-LDF with the ALM update) on OctMNIST L30/L40 $\times$ \{MNV3, RegNet, ViT\} $\times$ 4 seeds,
adjudicated against the frozen TraLO and Fioretto cells. \textbf{Found} (Table~\ref{tab:b3}): ALM is a
\emph{stronger} dual than plain Fioretto (it beats it by $+0.010$), yet TraLO tops ALM in all $6$ tight-cap
cells ($+0.028$ cc-F1). ALM does not close the gap, so the paper's justification for not adopting ALM stands.
This is the strongest new result.

\begin{table}[h]\centering\small
\begin{tabular}{lrcccc}
\toprule
comparison (cc-F1) & mean & W/T/L & MNV3 (L30/L40) & RegNet (L30/L40) & ViT (L30/L40) \\
\midrule
TraLO $-$ ALM (best trained base.) & $+0.028$ & 6/0/0 & $+0.009$ / $+0.013$ & $+0.048$ / $+0.008$ & $+0.046$ / $+0.046$ \\
TraLO $-$ Fioretto & $+0.039$ & 6/0/0 & $+0.016$ / $+0.022$ & $+0.035$ / $+0.021$ & $+0.086$ / $+0.051$ \\
ALM $-$ Fioretto & $+0.010$ & 5/0/1 & $+0.007$ / $+0.009$ & $-0.013$ / $+0.013$ & $+0.041$ / $+0.005$ \\
\bottomrule
\end{tabular}
\caption{B3: OctMNIST cc-F1 gaps, mean over 4 seeds; $6$ cells each. Raw cc-F1: TraLO $0.455 >$ ALM $0.427 >$
Fioretto $0.417$.}
\label{tab:b3}
\end{table}

\subsection{B4 -- Tie-region component ablation}
\textbf{Question.} Are the TraLO controls (reset / hinge / $\rho$-schedule / freeze) load-bearing in a
\emph{tie} region, or only at tight caps? \textbf{What ran.} $20$ runs: DermMNIST L50 (a tie cell), MNV3,
leave-one-out of each component $\times$ 4 seeds. \textbf{Found} (Table~\ref{tab:b4}): removing any single
component moves results by at most $\pm 0.006$ -- none is load-bearing in a tie region. This is the tighter
(and favourable) conclusion: the portable components carry the \emph{tight-cap} advantage specifically,
consistent with the mechanism.

\begin{table}[h]\centering\small
\begin{tabular}{lrr@{\hskip 2cm}lrr}
\toprule
removed & macro-F1 $\Delta$ & cc-F1 $\Delta$ & removed & macro-F1 $\Delta$ & cc-F1 $\Delta$ \\
\midrule
$-$ reset & $+0.002$ & $+0.006$ & $-$ $\rho$-schedule & $+0.002$ & $+0.003$ \\
$-$ hinge & $-0.003$ & $\phantom{+}0.000$ & $-$ freeze & $-0.002$ & $+0.003$ \\
\bottomrule
\end{tabular}
\caption{B4: leave-one-out deltas (full $-$ variant) on the DermMNIST L50 tie cell, mean over 4 seeds. Full run:
$0.744$ macro-F1 / $0.564$ cc-F1. No component is load-bearing here.}
\label{tab:b4}
\end{table}

\subsection{B5 -- Win-bar sensitivity}
\textbf{Question.} Is the OctMNIST regime classification stable across win thresholds $\tau$? \textbf{What ran.}
No new runs; the OctMNIST cc-F1 cells were re-scored at $\tau\in\{.003,.005,.010\}$. \textbf{Found}
(Table~\ref{tab:b5}): the win vs the trained baselines at L30/L40 is stable at every threshold ($4/0/0$ vs
Fioretto even at the strict $\tau=.010$), while the comparison against post-hoc clip never yields a stable win
at any threshold -- the honest cc-F1 tie.

\begin{table}[h]\centering\small
\begin{tabular}{llccc@{\hskip 1.2cm}llccc}
\toprule
vs.\ trained base. & cap & .003 & .005 & .010 & vs.\ clip (separate) & cap & .003 & .005 & .010 \\
\midrule
TraLO $-$ Fioretto & L30 & 4/0/0 & 4/0/0 & 4/0/0 & TraLO $-$ clip & L40 & 3/0/1 & 2/1/1 & 1/2/1 \\
                   & L40 & 4/0/0 & 4/0/0 & 4/0/0 &                & L30 & 2/0/2 & 2/0/2 & 0/2/2 \\
TraLO $-$ ALM      & L30 & 3/0/0 & 3/0/0 & 2/1/0 &                & L10--20 & tie & tie & tie \\
\bottomrule
\end{tabular}
\caption{B5: W/T/L across OctMNIST backbone cells by cap and win threshold $\tau$. Left: vs the trained
baselines (duals). Right: vs post-hoc clip (separate category).}
\label{tab:b5}
\end{table}

\subsection{B6 -- Bootstrap confidence intervals; B7 -- BH-FDR}
\textbf{Question.} Do the headline cells survive $n=4$ seeds statistically? \textbf{What ran.} No new runs;
$20{,}000$-resample bootstrap CIs over the OctMNIST backbone cells per cap, plus BH-FDR. \textbf{Found}
(Table~\ref{tab:b6}): CIs against the trained baselines \emph{exclude} $0$ at L30/L40 (the previously-flagged
ViT L30 cell is $+0.086$ vs Fioretto and survives), while the CI against post-hoc clip \emph{includes} $0$ at
every cap. The win is a mid-cap (L30/L40) phenomenon: at L10/L15/L20 even the dual gap ties because the
tightest caps ceiling-crush every method. \textbf{B7 (BH-FDR, $q=0.05$):} $3/8$ cells survive for
TraLO $-$ Fioretto; $0/8$ for TraLO $-$ clip.

\begin{table}[h]\centering\small
\begin{tabular}{llrl@{\hskip 1.4cm}llrl}
\toprule
vs.\ trained & cap & mean & 95\% CI & vs.\ clip (separate) & cap & mean & 95\% CI \\
\midrule
TraLO $-$ Fioretto & L30 & $+0.046$ & $[+.024,+.074]$ & TraLO $-$ clip & L10 & $+0.001$ & $[-.005,+.007]$ \\
                   & L40 & $+0.033$ & $[+.021,+.044]$ &                & L15 & $+0.004$ & $[\ .000,+.007]$ \\
                   & L20 & $+0.002$ & $[\ .000,+.005]$ &                & L20 & $+0.001$ & $[-.008,+.007]$ \\
TraLO $-$ ALM      & L30 & $+0.035$ & $[+.010,+.048]$ &                & L30 & $-0.003$ & $[-.013,+.008]$ \\
                   & L40 & $+0.022$ & $[+.008,+.046]$ &                & L40 & $+0.002$ & $[-.010,+.012]$ \\
\bottomrule
\end{tabular}
\caption{B6: grand-mean OctMNIST cc-F1 gap with $95\%$ bootstrap CI (resampling backbone cells). Trained-baseline
CIs exclude $0$ at L30/L40; the clip CI includes $0$ at every cap.}
\label{tab:b6}
\end{table}

\subsection{B8 -- Fourth backbone (MobileNetV2)}
\textbf{Question.} Does the cc-F1 win reproduce on a fourth backbone? \textbf{What ran.} $32$ runs: MobileNetV2
$\times$ OctMNIST L30/L40 $\times$ \{TraLO, Fioretto, Hounie, clip\} $\times$ 4 seeds. \textbf{Found}
(Table~\ref{tab:b8}): TraLO beats the best trained baseline (Fioretto) by $+0.042$ and Hounie collapses; the
comparison vs post-hoc clip is a thin cc-F1 tie ($+0.007$) with $0$ native satisfaction. The dual win is thus
backbone-general across MNV3/RegNet/ViT/MNV2.

\begin{table}[h]\centering\small
\begin{tabular}{lrrr@{\hskip 1.4cm}l}
\toprule
comparison (cc-F1) & mean & L30 & L40 & raw cc-F1 (MNV2) \\
\midrule
TraLO $-$ Fioretto (best trained base.) & $+0.042$ & $+0.047$ & $+0.037$ & TraLO $0.454$ \\
TraLO $-$ Hounie (collapses) & $+0.169$ & $+0.189$ & $+0.149$ & Fioretto $0.412$ \\
TraLO $-$ clip (separate category) & $+0.007$ & $+0.009$ & $+0.005$ & Hounie $0.285$, clip $0.447$ \\
\bottomrule
\end{tabular}
\caption{B8: MobileNetV2, OctMNIST cc-F1 gaps, mean over 4 seeds. The dual win reproduces on a fourth backbone.}
\label{tab:b8}
\end{table}

% =====================================================================
\section{What worked, what did not, and what was researched}
% =====================================================================
\textbf{What worked (goes into the paper).}
\begin{itemize}
\item TraLO beats the \emph{best trained baseline} on OctMNIST tight caps: $+0.028$ cc-F1 vs the strong ALM and
$+0.039$ vs Fioretto ($6/6$ cells), with bootstrap CIs excluding $0$ at L30/L40 (B3, B6) and stable under the
win-bar sweep and BH-FDR (B5, B7).
\item The win is \textbf{backbone-general}: it reproduces on MNV3, RegNet, ViT, and MNV2 (B8).
\item The controls are \textbf{mechanism-consistent}: they are load-bearing at tight caps but inert in a tie
region (B4), which is exactly what the mechanism predicts.
\item TraLO uniquely provides \textbf{native constraint satisfaction} ($1.0$) across every setting.
\end{itemize}
\textbf{What did not work / honest negatives (reported as scope).}
\begin{itemize}
\item Against \textbf{post-hoc clipping}, TraLO's cc-F1 is a \emph{tie}, not a win: the gap CI includes $0$ at
every cap (B5, B6). Clipping matches accuracy but does not satisfy the cap natively, so the honest claim vs
clip is ``parity $+$ native satisfaction,'' not an accuracy win.
\item At \textbf{native resolution} the tight-cap advantage does \emph{not} appear: TraLO ties the best trained
baseline on all five datasets (B2). An earlier apparent TissueNative edge did not survive adding a backbone and
a tighter cap.
\item \textbf{Focal loss} matches or beats TraLO's macro-F1 on DermMNIST/MNV3 (B1); the macro-F1 story is
backbone-dependent, not a clean TraLO win.
\end{itemize}
\textbf{What was researched (the scientific contribution).} The \emph{headroom mechanism} (Section~7) --- a
precise characterization of \emph{when} transductive count-constraint training helps --- which unifies every
win and every tie above.

% =====================================================================
\section{The mechanism}
% =====================================================================
The count penalty (\ref{eq:pen}) reaches the network weights only through the scalar soft-count $S_{c^\star}$,
so its gradient is a \emph{scalar $\times$ a fixed direction} and its overall scale ($\lambda,\rho$) is absorbed
by the Adam optimizer and gradient clipping. What the penalty can \emph{do} therefore depends entirely on
whether CE is still active:
\begin{itemize}
\item \textbf{CE active} (tight cap $+$ warm-up headroom): the penalty \emph{re-ranks} borderline examples,
producing a genuine cc-F1 gain over the trained duals (B3, B8).
\item \textbf{CE saturated} (loose cap, long warm-up, easy data, or native resolution): the penalty can only
\emph{shift} logits uniformly, leaving the top-$K$ set --- and thus the metric --- unchanged, giving a tie
(B2, B4, B5--B6).
\end{itemize}
Post-hoc clipping sits on a separate axis: it can match cc-F1 but, being applied after training, can never
satisfy the cap natively. This single mechanism predicts every Track~B win and tie, and is supported by the
Adam scale-invariance analysis in the main text.

% =====================================================================
\section{Limitations and honest scope}
% =====================================================================
The cc-F1 advantage is over the trained \emph{duals}, and only at mid tight caps (L30/L40); it narrows to a tie
at the very tightest caps (L10--L20), where every method is ceiling-crushed. Against post-hoc clipping it is a
tie plus native satisfaction. Focal ties TraLO on macro-F1 on DermMNIST. Native resolution is a
trained-baseline tie. cc-F1 gaps are small by construction (ceiling $2c/(1+c)$). Framed as a characterization
of \emph{when} count-constraint training helps --- beating the trained duals where the cap binds and the network
is unsaturated, matching strong baselines elsewhere, and uniquely satisfying the cap natively --- this is a
precise and defensible contribution.

\end{document}
